%=====================================================================
\chapter{Snowballing, Scoping Reviews and Grey Literature}\label{ch:greylit}
%=====================================================================
\locator{2}{Part~2 --- Finding and Synthesising Knowledge}

\begin{outcomes}
By the end of this chapter you will be able to:
\begin{tight}
  \item \textbf{Conduct} backward and forward snowballing from a defensible
    start set, and state a stopping rule.
  \item \textbf{Choose} between a systematic review, a scoping review, a rapid
    review and a multivocal review.
  \item \textbf{Justify} the inclusion of grey literature, and appraise it with
    criteria designed for it.
  \item \textbf{Explain} how publication bias distorts a database-only search,
    and what partially corrects it.
\end{tight}
\end{outcomes}

\begin{prereq}
\chref{slr}. This chapter is the set of techniques that repair the database
search, and the set of review types that are appropriate when a full systematic
review is not.
\end{prereq}

\begin{keyterms}
snowballing $\bullet$ backward and forward snowballing $\bullet$ start set
$\bullet$ scoping review $\bullet$ rapid review $\bullet$ grey literature
$\bullet$ multivocal literature review $\bullet$ publication bias $\bullet$
file drawer problem $\bullet$ saturation of search
\end{keyterms}

\section{Why the Database Search Is Never Enough}

A Boolean string matches terminology, and terminology is unstable. Authors in
adjacent communities describe the same phenomenon differently; a foundational
paper may predate the term now used for its subject; a highly relevant study
may use ``serverless'' where you searched for ``function-as-a-service''. No
string catches all of this, and a review that relies on one alone has a
coverage problem it cannot see.

Snowballing exploits a different signal: what authors themselves considered
relevant.

\begin{definitionbox}[Backward and forward snowballing]
\term{Backward snowballing} follows the reference list of an included paper to
find earlier work it drew on. \term{Forward snowballing} finds later papers
that cite it, using Google Scholar, Scopus or Semantic Scholar.

\medskip
Together they traverse the citation graph in both directions from a
\term{start set} of known-relevant papers, iterating until no new relevant
papers appear.
\end{definitionbox}

\begin{paperbox}[Snowballing as a systematic procedure]
Claes Wohlin, \emph{Guidelines for snowballing in systematic literature studies
and a replication in software engineering}, EASE
2014~\cite{wohlin2014guidelines}.

\medskip
\textbf{Why it matters.} It makes snowballing a documented procedure rather
than incidental reading, and shows it can produce results comparable to a
database search --- sometimes better, because it is insensitive to terminology
drift.

\medskip
\textbf{What to extract.} How to choose a start set that is diverse in
community, venue and year; the iteration procedure; and the discipline of
recording, at each iteration, how many papers were examined and how many
included.
\end{paperbox}

\subsection{The procedure}

\begin{center}
\small
\begin{tabular}{@{}L{1.4cm}L{13.0cm}@{}}
\toprule
Step 1 & \textbf{Choose the start set.} Diverse in authorship, venue,
         community and publication year. A start set of five papers by one
         group traverses only that group's citation neighbourhood. \\
Step 2 & \textbf{Backward pass.} For each paper, scan the reference list.
         Screen on title, then on the \emph{context of the citation} --- read
         the sentence that cites it. This is faster and more accurate than
         reading abstracts. \\
Step 3 & \textbf{Forward pass.} For each paper, list citing works and screen
         the same way. \\
Step 4 & \textbf{Iterate.} New included papers form the next iteration's set. \\
Step 5 & \textbf{Stop} when an iteration yields no new included papers. Record
         the iteration count and the yield at each. \\
\bottomrule
\end{tabular}
\end{center}

\begin{conceptbox}[Reading the citation context is the trick]
The single most useful habit in snowballing: do not screen a reference by its
title. Screen it by \emph{what the citing paper says about it}. ``Similar to
the approach of Ahmed et al.\ [12], we partition\ldots'' tells you far more
about the relevance of [12] than its title does, and takes two seconds.
\end{conceptbox}

\section{Four Kinds of Review, and When Each Is Right}

\begin{center}
\small
\begin{longtable}{@{}L{2.5cm}L{4.3cm}L{3.4cm}L{4.0cm}@{}}
\toprule
\textbf{Type} & \textbf{Question it answers} & \textbf{Effort} & \textbf{Use when} \\
\midrule
\endfirsthead
\toprule
\textbf{Type} & \textbf{Question it answers} & \textbf{Effort} & \textbf{Use when} \\
\midrule
\endhead
\bottomrule
\endfoot
Systematic review
  & What is the effect of X on Y?
  & Months; two screeners
  & The question is specific and enough primary studies exist \\
Systematic mapping
  & What has been researched, and where are the gaps?
  & Weeks to months
  & Entering an area; generating the question a review will
    answer~\cite{petersen2015guidelines} \\
Scoping review
  & What is the extent and nature of the evidence?
  & Weeks
  & The concept itself is unclear and needs charting before it can be
    questioned \\
Rapid review
  & What does the evidence say, quickly?
  & Days to weeks
  & A decision must be made on a deadline; limitations declared explicitly \\
Multivocal (MLR)
  & What do researchers \emph{and} practitioners know?
  & As a review, plus grey literature appraisal
  & The topic is practice-led and the academic literature lags it \\
\end{longtable}
\end{center}

\begin{alertbox}[A rapid review is legitimate; a sloppy review called rapid is not]
Rapid reviews shorten the process deliberately and say exactly where. Acceptable
shortcuts: fewer databases, a restricted date range, single screening with a
sampled check, no formal quality scoring. What is never acceptable, in any
review type: an undocumented search, criteria invented during screening, or a
synthesis that does not report what was excluded. Speed is bought from
coverage, not from transparency.
\end{alertbox}

\section{Grey Literature}

\term{Grey literature} is everything produced outside commercial and academic
publishing: technical reports, standards, white papers, engineering blogs,
conference talks, documentation, Stack Overflow, theses, and preprints.

For large parts of computing this is not a marginal supplement. Industrial
practice frequently runs years ahead of the academic literature --- container
orchestration, site reliability engineering, and much of modern MLOps were
documented in engineering blogs and internal reports long before they were
studied. A review of such a topic that excludes grey literature is not
conservative; it is describing a different, later world.

\begin{paperbox}[When and how to include the practitioner literature]
Vahid Garousi, Michael Felderer and Mika M{\"a}ntyl{\"a}, \emph{Guidelines for
including grey literature and conducting multivocal literature reviews in
software engineering}, Information and Software Technology 106,
2019~\cite{garousi2019guidelines}.

\medskip
\textbf{Why it matters.} It supplies both halves of the problem: a decision
procedure for whether grey literature is needed at all, and a quality
instrument designed for sources that have no peer review to lean on.

\medskip
\textbf{What to extract.} Their seven questions for deciding to include grey
literature --- of which the decisive ones are whether the topic is
practice-led, whether the academic literature is sparse, and whether
practitioner voice is itself part of the phenomenon --- and their tiered
classification of grey sources by expertise and outlet control.
\end{paperbox}

\subsection{Appraising a source with no peer review}

Peer-reviewed papers arrive with a weak quality signal already attached. Grey
literature does not, so appraisal has to do more work.

\begin{center}
\small
\begin{tabular}{@{}L{2.7cm}L{11.6cm}@{}}
\toprule
\textbf{Criterion} & \textbf{Question to ask} \\
\midrule
Authority   & Is the author identifiable, and do they have demonstrable
              standing --- a role, a track record, an organisation? \\
Methodology & Is it stated how the claims were arrived at? A post reporting
              measured latencies with a described setup differs entirely from
              one asserting a conclusion \\
Objectivity & Is there a commercial interest? A vendor blog showing its own
              product winning a benchmark is evidence about the vendor's
              claims, not about the product \\
Date        & Is it dated, and is it current for a fast-moving topic? \\
Novelty     & Does it add something the peer-reviewed sources do not? \\
Impact      & Is it cited, referenced or discussed by others in the field? \\
Outlet type & Tier 1: books, reports, theses. Tier 2: annual reviews,
              presentations, wikis. Tier 3: blogs, emails, tweets \\
\bottomrule
\end{tabular}
\end{center}

\begin{pitfall}[Three ways grey literature is misused]
\textbf{Cited as if it were a study.} ``Blogs report that microservices improve
scalability [17]'' treats an opinion as a finding. Attribute it: ``a
practitioner account describes\ldots''. The claim is about what practitioners
believe, which may be exactly what you want --- but say so.

\medskip
\textbf{Included without appraisal.} If peer-reviewed studies are quality-scored
and blog posts are not, the grey literature has been given an easier standard
than the academic literature, which is the opposite of the correct treatment.

\medskip
\textbf{Uncitable because it moved.} Blogs vanish and get edited. Archive every
grey source --- the Internet Archive's Wayback Machine will do --- and cite the
archived URL with the access date. A reference an examiner cannot follow is not
a reference.
\end{pitfall}

\section{Publication Bias}

Studies finding an effect are more likely to be submitted, more likely to be
accepted, and published sooner than studies finding nothing. The consequence is
that the literature you can find is not the research that was done, and a
perfectly executed review of a biased literature returns a biased answer with
false confidence. The unpublished null results are the \term{file drawer}.

In medicine this is partly controlled by trial registries: a study registered
before it runs can be chased even if never published. Computing has no
equivalent, which makes the problem harder for us, not easier.

\begin{conceptbox}[What you can actually do about it]
\begin{tightnum}
  \item \textbf{Search for the unpublished.} Theses, technical reports and
    preprints are where negative results in computing usually end up. This is
    a strong argument for grey literature even in a review that is otherwise
    academic.
  \item \textbf{Look at the shape of the evidence.} If every small study finds
    a large effect and the only large study finds none, that pattern is what
    publication bias looks like. Where enough effect sizes exist, a funnel
    plot makes it visible.
  \item \textbf{Say so.} A review that names publication bias as a limitation
    and explains which direction it would push the conclusion is more useful
    than one that ignores it.
  \item \textbf{Do not add to it.} Pre-register and report your own null
    results (\chref{prereg}).
\end{tightnum}
\end{conceptbox}

\begin{traditions}{extending the search}
\tradEMP{Snowballing improves coverage of comparable studies; grey literature
matters mainly as a source of unpublished nulls and of effect sizes omitted
from published abstracts.}
\tradDSR{Practitioner sources often \emph{are} the state of the art. Design
knowledge frequently appears first in engineering blogs and standards, and a
DSR review that ignores them misdescribes the problem space.}
\tradQUAL{Practitioner writing is data about how a community understands its
own work, and can be analysed as documents rather than merely cited as
evidence.}
\tradFORM{Snowballing is the primary method: results are connected by
citation far more reliably than by keyword, and terminology in theory papers
is comparatively stable.}
\end{traditions}

\begin{lab}[Snowball from three seeds]
Take three papers from your quasi-gold standard in \chref{slr}, chosen to be
from different groups and different years.

\begin{tightnum}
  \item Backward pass: screen every reference by citation context. Record
    examined and included counts.
  \item Forward pass: use Google Scholar or Semantic Scholar. Record the same.
  \item Run a second iteration on whatever you included.
  \item Compare the result with your database search. How many papers did
    snowballing find that the string missed? How many did the string find that
    snowballing missed?
\end{tightnum}

\medskip
The overlap tells you how good your string was. A large snowballing-only set
means the string needs repair --- go back and repair it, then report both in
the PRISMA flow.
\end{lab}

\begin{checkpoint}
\begin{tightnum}
  \item Why is a start set of five papers by the same research group a poor
    one, even if all five are highly relevant?
  \item Give a computing topic for which excluding grey literature would
    produce a misleading review, and say why.
  \item A vendor benchmark shows their database is fastest. Under what
    circumstances, if any, would you include it in a review, and how would you
    describe it?
  \item Publication bias inflates apparent effects. Name a situation in
    computing research where it could plausibly \emph{deflate} them instead.
\end{tightnum}
\end{checkpoint}

\begin{chaptersummary}
\begin{tight}
  \item Boolean search matches terminology; snowballing follows what authors
    themselves judged relevant. Use both, and report the yield of each.
  \item Screen references by citation context, not by title.
  \item Mapping, scoping, rapid and multivocal reviews are legitimate types
    with different purposes. Speed is bought from coverage, never from
    transparency.
  \item In practice-led computing topics grey literature is not a supplement;
    excluding it describes a different world.
  \item Appraise grey sources on authority, methodology, objectivity, date,
    novelty, impact and outlet tier --- and archive them, or they will vanish.
  \item Publication bias means the findable literature is not the research
    done. Search for theses and preprints, look at the shape of the evidence,
    declare the limitation, and do not add to the problem yourself.
\end{tight}
\end{chaptersummary}

\begin{reviewq}
  \item Wohlin reports snowballing producing results comparable to database
    search. If that is so, why do guidelines still require the database search?
    Give the strongest argument for keeping both.
  \item Grey literature inclusion is contested: critics say it imports
    unreviewed claims into the evidence base. Argue the case against inclusion
    as strongly as you can, then say what would answer it.
  \item Computing has no trial registry. Design one that could plausibly be
    adopted by a software engineering conference, and identify the incentive
    problem that would most likely defeat it.
  \item A multivocal review finds that practitioners and researchers disagree
    about whether a technique works. Is that a limitation of the review, or a
    finding? Defend your answer.
\end{reviewq}
